% Generated draft: requires mathematical and source review.
\documentclass[11pt]{article}
\usepackage[margin=25mm]{geometry}
\usepackage{fontspec}
\setmainfont{Latin Modern Roman}
\usepackage{amsmath,amssymb,mathtools,unicode-math}
\setmathfont{Latin Modern Math}
\usepackage{enumitem,needspace,xurl,hyperref}
\setlength{\parindent}{0pt}
\setlength{\parskip}{5pt}
\setlength{\emergencystretch}{3em}
\newcommand{\N}{\mathbb{N}}
\newcommand{\Z}{\mathbb{Z}}
\newcommand{\Q}{\mathbb{Q}}
\newcommand{\R}{\mathbb{R}}
\newcommand{\C}{\mathbb{C}}
\newcommand{\F}{\mathbb{F}}
\newcommand{\A}{\mathbb{A}}
\newcommand{\PP}{\mathbb{P}}
\newcommand{\E}{\mathbb{E}}
\newcommand{\eps}{\varepsilon}
\newcommand{\dd}{\,\mathrm{d}}
\begin{document}
\section*{UnsolvedMath ID: 4000007}
\textbf{Source problem code:} catalogue entry
\subsection{Definitions and mathematical statement}
Let \((X,d)\) be a metric space equipped with a Markov kernel \(P\), and write \(P_x=P(x,\cdot)\). Assume the probability laws under consideration have finite first moments. For \(x\ne y\), define
\[
\kappa_1(x,y)=1-\frac{W_1(P_x,P_y)}{d(x,y)},\qquad
\kappa_t(x,y)=1-\frac{W_1(P^t(x,\cdot),P^t(y,\cdot))}{d(x,y)}
\quad(t\in\N_{>0}),
\]
where \(P^t\) is the iterated kernel and \(W_1\) is the infimum of \(\int d(u,v)\,\mathrm d\pi(u,v)\) over couplings of its two arguments.

Ollivier's Problem G asks for two extensions of uniform-curvature estimates. First, can such estimates survive when curvature is nonpositive only at a controlled collection of locations? At the pair-curvature level the exceptional set would be
\[
B=\{(x,y)\in X^2:x\ne y,\ \kappa_1(x,y)\le0\}.
\]
The source does not say whether the control is on this set, on graph edges, or on a pointwise curvature derived from it, nor whether number means cardinality, probability mass, or another quantitative defect.

Second, can \(\kappa_t\) be quantitatively related to curvature averaged along random-walk trajectories, perhaps using large deviations? For a chain \((X_i)_{i\ge0}\) with kernel \(P\), a trajectory average would have the form
\[
A_t=\frac1t\sum_{i=0}^{t-1}c(X_i,X_{i+1}),
\]
where \(c\) is the chosen curvature observable. The source supplies neither that observable nor a comparison inequality, tail event, deviation rate, starting-law convention, or hypotheses ensuring such a result; the formula therefore specifies the kind of average, not an asserted theorem. In particular, pair curvature is undefined on diagonal steps unless a convention is supplied.

The supplied research summary reports partial progress through averaged or decaying curvature diameter and comparison results, but explicitly does not report a fully general trajectory-average large-deviation theorem. These more specialized results are not treated as a completion of the two-part programme.
\subsection{Short English statement}
Extend coarse-Ricci estimates beyond uniform lower bounds, allowing controlled nonpositive curvature, and relate iterated-kernel curvature to trajectory averages. The source leaves the defect control and proposed large-deviation statement unspecified.
\subsection{Sources}
\begin{itemize}
\item Local UnsolvedMath registry, permanent ID \(4000007\), exact source code \texttt{catalogue entry}, ``Non-constant curvature.''
\item Yann Ollivier, \emph{Discrete Ricci curvature: Open problems} (2008), Problem G; \url{https://www.yann-ollivier.org/rech/publs/problems_curvmarkov.pdf}. Recorded access 2026-07-29; pdf-text extraction.
\item Supplied research summary refers to Ollivier's main article and to Münch, \emph{Advances in Mathematics} (2019), for a diameter/Laplacian-comparison line involving curvature decay and average conditions. Complete publication details and exact theorem statements are not supplied.
\end{itemize}
\end{document}
